% TOP_PROBLEM: {"id": 20000950, "title": "A small Fano-good tight path and the connectedness convention", "category_id": 2, "category": {"id": 2, "name": "combinatorics", "display_name": "Combinatorics", "description": "Counting problems, graph theory, discrete structures.", "slug": "combinatorics", "order_index": 2, "created_at": "2026-07-31T15:26:25.671Z"}, "status": "partially_solved", "problem_number": "AIM-COMBINATORICS-0075", "set_id": 14, "statusQualification": "Distinguishes the literal equality defining Fano-goodness from the connectedness assumption needed for the accompanying standard lower-bound construction."}
% FIRST_POSTED: 2026-10-01
% ENTRY_KIND: statement_only
% UnsolvedMath ID 20000950; source catalogue code AIM-COMBINATORICS-0075
% !TeX program = lualatex
% Definition and sources only; this file is not an LLM solution attempt.
\documentclass[11pt,a4paper]{article}
\usepackage[margin=25mm]{geometry}
\usepackage{fontspec}
\setmainfont{lmroman10-regular.otf}[BoldFont=lmroman10-bold.otf,
 ItalicFont=lmroman10-italic.otf,BoldItalicFont=lmroman10-bolditalic.otf]
\usepackage{amsmath,amssymb,mathtools}
\usepackage{unicode-math}
\setmathfont{latinmodern-math.otf}
\AtBeginDocument{\let\setminus\smallsetminus}
\usepackage{xurl,hyperref}
\hypersetup{colorlinks=true,urlcolor=blue}
\setcounter{secnumdepth}{0}
\setlength{\parindent}{0pt}
\setlength{\parskip}{5pt}
\setlength{\emergencystretch}{3em}
\begin{document}
\section*{A small Fano-good tight path and the connectedness convention}\label{20000950}
{\small UnsolvedMath ID 20000950 · AIM-COMBINATORICS-0075 · Combinatorics · AIM Workshop Problem Lists}
\subsection{Definitions and mathematical statement}
A finite simple $3$-uniform hypergraph consists of a vertex set and a set of three-element subsets, called edges.
Let $F$ be the Fano plane: its vertices are the seven nonzero vectors of $\mathbb F_2^3$, and its edges are the unordered triples $\{x,y,x+y\}$ with $x\ne y$ nonzero.
These give seven distinct edges.
For finite $3$-uniform hypergraphs $H,J$, let $r(H,J)$ be the least $N$ such that every red--blue coloring of the triples of an $N$-element set contains a red copy of $H$ or a blue copy of $J$.
A copy is an injective vertex map carrying every required edge to an edge of the designated color; induced containment is not required.

Classify the hypergraphs $H$ for which
\[
 r(H,F)=2\bigl(|V(H)|-1\bigr)+1.
\]
This equality is the workshop's definition of being Fano-good.
The question may be considered for all finite $H$ having at least one edge, including the role of isolated vertices.
For connected $H$, the usual lower-bound coloring divides $2(|V(H)|-1)$ vertices into two equal parts, colors triples within a part red, and mixed triples blue.
Here connected means connectedness of the vertex--edge incidence graph.
That construction's justification uses connectedness of $H$, whereas the numerical equality itself remains a meaningful classification question without it.
A solution should identify exactly which structural properties force equality rather than a different Ramsey threshold.
\subsection{Short English statement}
Which three-uniform hypergraphs attain the proposed exact Ramsey threshold against the Fano plane?
\subsection{Sources}
\begin{itemize}
\item[S1] ULAM.AI and the UnsolvedMath Contributors, supplied problems.json, record 20000950 (AIM-COMBINATORICS-0075), AIM Workshop Problem Lists. Catalogue identity and source transcription; CC BY 4.0.
\url{https://huggingface.co/datasets/ulamai/UnsolvedMath}.
\item[S2] American Institute of Mathematics, Graph Ramsey theory, item 28. Original workshop question underlying AIM-COMBINATORICS-0075.
\url{https://www.overleaf.com/read/mnvcscjjysvg}.
\end{itemize}
\end{document}
